\begin{name}
	{ÔN TẬP KIỂM TRA CUỐI HỌC KÌ 1 - KNTT}
	{TOÁN 10}
	{LỚP TOÁN THẦY PHÁT}
	{\thoigian}
\end{name}

%================================PHẦN I
\caulc
\Opensolutionfile{ans}[ans/0-HK1-KN-4-ThanhBinh-DongNai-2324-Phan-I]

%%%==============EX_1============%%%
\begin{ex}%[0-HK1-KN-4-ThanhBinh-DongNa-2324]%[VN-MT-25, Phạm Ngọc Trung]%[0D1N1-1]
Phát biểu nào sau đây là một mệnh đề?
\choice
{$2$K$8$ ơi, cố gắng nhé!}
{\True Số $2024$ là số tự nhiên chẵn}
{Đề toán hôm nay có dễ không?}
{Chúc các em làm bài tốt!}
\loigiai{\lq\lq Số $2024$ là số tự nhiên chẵn\rq\rq\, là một khẳng định đúng nên là mệnh đề.}
\end{ex}
%%%==============EX_2============%%%
\begin{ex}%[0-HK1-KN-4-ThanhBinh-DongNa-2324]%[VN-MT-25, Phạm Ngọc Trung]%[0D2N1-1]
Bất phương trình nào sau đây \textbf{không} phải là bất phương trình bậc nhất hai ẩn?
\choice
{$2x-3y<0$}
{\True $x^2+y\leq 5$}
{$-x+5y\geq-2$}
{$x+y>-6$}
\loigiai{$x^2+y\leq 5$ không phải là bất phương trình bậc nhất hai ẩn vì chứa $x^2$ không phải bậc nhất.}
\end{ex}
%%%==============EX_3============%%%
\begin{ex}%[0-HK1-KN-4-ThanhBinh-DongNa-2324]%[VN-MT-25, Phạm Ngọc Trung]%[0H4N2-1]
Xét tam giác $ABC$ tùy ý có $BC=a$, $AC=b$, $AB=c$. Mệnh đề nào dưới đây đúng?
\choice
{$b^2=a^2+c^2+2ac\cos B$}
{$b^2=a^2+c^2-ac\cos B$}
{\True $b^2=a^2+c^2-2ac\cos B$}
{$b^2=a^2+c^2+ac\cos B$}
\loigiai{Định lý cosin $b^2=a^2+c^2-2ac \cos B$.}
\end{ex}
%%%==============EX_4============%%%
\begin{ex}%[0-HK1-KN-4-ThanhBinh-DongNa-2324]%[VN-MT-25, Phạm Ngọc Trung]%[0D6N1-2]
Độ cao của một ngọn núi được ghi lại như sau $\bar{h}=1372{,}5\mathrm{~m} \pm 0{,}3\mathrm{~m}$. Độ chính xác $d$ của phép đo trên là
\choice
{$d=0{,}1$ m}
{$d=\pm 0{,}1$ m}
{$d=\pm 0{,}3$ m}
{\True $d=0{,}3$ m}
\loigiai{Độ chính xác $d$ của phép đo trên là $d=0{,}3$ m.}
\end{ex}
%%%==============EX_5============%%%
\begin{ex}%[0-HK1-KN-4-ThanhBinh-DongNa-2324]%[VN-MT-25, Phạm Ngọc Trung]%[0D6N3-1]
Số liệu đặc trưng nào sau đây đo xu thế trung tâm của mẫu số liệu thông kê?
\choice
{Khoảng biến thiên}
{\True Số trung bình}
{Phương sai}
{Độ lệch chuẩn}
\loigiai{Số liệu đặc trưng đo xu thế trung tâm của mẫu số liệu thông kê là số trung bình.}
\end{ex}
%%%==============EX_6============%%%
\begin{ex}%[0-HK1-KN-4-ThanhBinh-DongNa-2324]%[VN-MT-25, Phạm Ngọc Trung]%[0D6N4-1]
Khoảng biến thiên của mẫu số liệu là
\choice
{hiệu số giữa tứ phân vị thứ ba và tứ phân vị thứ nhất}
{\True hiệu số giữa số lớn nhất và số nhỏ nhất trong mẫu số liệu đó}
{căn bậc hai của phương sai}
{bình phương của phương sai}
\loigiai{Khoảng biến thiên của mẫu số liệu là hiệu số giữa số lớn nhất và số nhỏ nhất trong mẫu số liệu đó.}
\end{ex}
%%%==============EX_7============%%%
\begin{ex}%[0-HK1-KN-4-ThanhBinh-DongNa-2324]%[VN-MT-25, Phạm Ngọc Trung]%[0D6N4-4]
Mẫu số liệu sau đây cho biết sĩ số của $5$ lớp khối $10$ tại một trường
\begin{center}
\begin{tabular}{lllll}
	$43$&$45$& $44$& $38$& $40$\\
\end{tabular}
\end{center}
Biết phương sai của mẫu số liệu này là $6{,}8$. Độ lệch chuẩn của mẫu số liệu này là
\choice
{$13{,}6$}
{\True $\sqrt{6,8}$}
{$(6{,}8)^2$}
{$3{,}4$}
\loigiai{Độ lệch chuẩn của mẫu số liệu bằng căn bậc hai của phương sai nghĩa là bằng $\sqrt{6,8}$.}
\end{ex}
%%%==============EX_8============%%%
\begin{ex}%[0-HK1-KN-4-ThanhBinh-DongNa-2324]%[VN-MT-25, Phạm Ngọc Trung]%[0H5N1-1]
Cho hình bình hành $ABCD$, vectơ có điểm đầu và điểm cuối là các đỉnh của hình bình hành bằng với vectơ $\overrightarrow{DC}$ là
\choice
{\True $\overrightarrow{AB}$}
{$\overrightarrow{BA}$}
{$\overrightarrow{CD}$}
{$\overrightarrow{AC}$}
\loigiai{$\overrightarrow{DC}=\overrightarrow{AB}$ vì hai vectơ này cùng hướng và cùng độ dài.}
\end{ex}
%%%==============EX_9============%%%
\begin{ex}%[0-HK1-KN-4-ThanhBinh-DongNa-2324]%[VN-MT-25, Phạm Ngọc Trung]%[0H5N2-1]
Cho ba điểm $M$, $N$, $P$ bất kì. Khi đó $\overrightarrow{PN}-\overrightarrow{PM}$ bằng
\choice
{\True $\overrightarrow{MP}$}
{$\overrightarrow{PM}$}
{$\overrightarrow{MN}$}
{$\overrightarrow{NM}$}
\loigiai{Theo quy tắc hiệu của hai vectơ ta có $\overrightarrow{PN}-\overrightarrow{PM}=\overrightarrow{MN}$.}
\end{ex}
%%%==============EX_10============%%%
\begin{ex}%[0-HK1-KN-4-ThanhBinh-DongNa-2324]%[VN-MT-25, Phạm Ngọc Trung]%[0H9N1-1]
Trong mặt phẳng tọa độ $Oxy$, cho hai điểm $M(4;-3)$, $N(-2; 0)$. Tọa độ của vectơ $\overrightarrow{MN}$ là
\choice
{$(4;-5)$}
{\True $(-6; 3)$}
{$(-4; 5)$}
{$(6;-3)$}
\loigiai{Từ định nghĩa suy ra $\overrightarrow{MN}=(-6; 3)$.}
\end{ex}
%%%==============EX_11============%%%
\begin{ex}%[0-HK1-KN-4-ThanhBinh-DongNa-2324]%[VN-MT-25, Phạm Ngọc Trung]%[0H9N1-2]
Trong mặt phẳng tọa độ $Oxy$, cho hai vectơ $\overrightarrow{a}=(1; 4)$ và $\overrightarrow{b}=(-3; 2)$.
Tính tích vô hướng $\overrightarrow{a} \cdot \overrightarrow{b}$
\choice
{$\overrightarrow{a} \cdot \overrightarrow{b}=-16$}
{\True $\overrightarrow{a} \cdot \overrightarrow{b}=5$}
{$\overrightarrow{a} \cdot \overrightarrow{b}=-5$}
{$\overrightarrow{a} \cdot \overrightarrow{b}=11$}
\loigiai{Ta có $\overrightarrow{a} \cdot \overrightarrow{b}=1\cdot(-3)+4\cdot 2=5$.
}
\end{ex}
%%%==============EX_12============%%%
\begin{ex}%[0-HK1-KN-4-ThanhBinh-DongNa-2324]%[VN-MT-25, Phạm Ngọc Trung]%[0D3N1-1]
Tập xác định của hàm số $y=\dfrac{1}{x-6}$ là
\choice
{$(6;+\infty)$}
{$[6;+\infty)$}
{$(-\infty;+\infty)$}
{\True $\mathbb{R} \setminus\{6\}$}
\loigiai{Hàm số xác định $\Leftrightarrow x-6\neq 0\Leftrightarrow x \neq 6$.\\
Vậy tập xác định của hàm số $y=\dfrac{1}{x-6}$ là $\mathbb{R} \setminus\{6\}$.
}
\end{ex}
\Closesolutionfile{ans}
%================================PHẦN II
\cauds
\Opensolutionfile{ans}[ans/0-HK1-KN-4-ThanhBinh-DongNai-2324-Phan-II]
%%%==============EX_1============%%%
\begin{ex}%[0-HK1-KN-4-ThanhBinh-DongNa-2324]%[VN-MT-25, Phạm Ngọc Trung]%[0H5H3-2]
	Cho tam giác $ABC$ cân tại $A$ có $G$ là trọng tâm, $I$ là trung điểm đoạn thẳng $BC$.
	\choiceTF
	{$\overrightarrow{AB}=\overrightarrow{AC}$}
	{$\overrightarrow{IB}=\overrightarrow{IC}$}
	{\True $\overrightarrow{GB}+\overrightarrow{GC}=2\overrightarrow{GI}$}
	{$\overrightarrow{MA}+\overrightarrow{MB}+\overrightarrow{MC}=-3\overrightarrow{MG}$}
	\loigiai{
		\begin{itemchoice}
			\itemch \textbf{Sai}\\
			Ta có $\overrightarrow{AB}$, $\overrightarrow{AC}$ là hai vectơ không cùng phương nên $\overrightarrow{AB}=\overrightarrow{AC}$ là khẳng định sai.
			\itemch \textbf{Sai}\\
			Ta có $\overrightarrow{IB}$, $\overrightarrow{IC}$ là hai vectơ ngược hướng nhau nên $\overrightarrow{IB}=\overrightarrow{IC}$ là khẳng định sai.
			\itemch \textbf{Đúng}\\
			Do $I$ là trung điểm $BC$ ta có $\overrightarrow{GB}+\overrightarrow{GC}=2\overrightarrow{GI}$.
			\itemch \textbf{Sai}\\
			Do $G$ là trọng tâm tam giác $ABC$ ta có $\overrightarrow{MA}+\overrightarrow{MB}+\overrightarrow{MC}=3\overrightarrow{MG}$.
		\end{itemchoice}
		
	}
\end{ex}

\begin{ex}%[0-HK1-KN-4-ThanhBinh-DongNa-2324]%[VN-MT-25, Phạm Ngọc Trung]%[0D3N2-4]
	Cho hàm số $y=f(x)=x^2$ có đồ thị $(C)$ và hàm số $y=g(x)=-2x+3$ có đồ thị là đường thẳng $(d)$.
	\choiceTF
	{\True Điểm $M(-1;1)$ thuộc đồ thị $(C)$}
	{Hàm số $y=f(x)$ nghịch biến trên $(1;5)$}
	{Đồ thị hàm số $(C)$ đi qua giao điểm của đường thẳng $(d)$ với trục tung}
	{\True Điểm $A(1;1)$ và $B(-3;9)$ là hai giao điểm của đồ thị $(C)$ và đường thẳng $(d)$}
	\loigiai{
		\begin{itemchoice}
			\itemch \textbf{Đúng}\\
			Ta có $f(-1)=1$ nên $M(-1; 1) \in(C)$.
			\itemch \textbf{Sai}\\
			Hàm số $y=f(x)=x^2$ đồng biến trên $(0;+\infty)$ và nghịch biến trên $(-\infty;0)$.\\
			Lại có $(1;5) \subset(0;+\infty)$ nên hàm số đồng biến trên $(1;5)$.
			\itemch	\textbf{Sai}\\
			Đường thẳng $(d)$ cắt trục tung tại điểm có $N(0;3)$. Ta có $f(0)=0$ nên điểm $N(0;3)$ không thuộc đồ thị hàm số $(C)$.
			\itemch \textbf{Đúng}\\
			Thay tọa độ điểm $A(1;1)$ và $B(-3;9)$ vào phương trình của $(C)$ và $(d)$ đều thỏa mãn.
		\end{itemchoice}
	}
\end{ex}

\begin{ex}%[0-HK1-KN-4-ThanhBinh-DongNa-2324]%[VN-MT-25, Phạm Ngọc Trung]%[0D1V2-3]
	Cho các tập hợp $A=(-\infty ; m)$, $B=[3m-1;3m+3]$ với $m \in \mathbb{R}$.
	\choiceTF
	{Khi $m=1$ thì tập $B$ chứa $3$ giá trị nguyên}
	{\True Với $m=2$ thì $A \cap B=\varnothing$}
	{Các giá trị của $m$ để $B \subset A$ là $m \geq \dfrac{1}{2}$}
	{\True Các giá trị của $m$ để $C_{\mathbb{R}}A \cap B \neq \varnothing$ là $m \geq-\dfrac{3}{2}$}
	\loigiai{
		\begin{center}
		\begin{tikzpicture}[line join = round, line cap = round,>=stealth,font=\footnotesize,scale=1]
			\def\xmin{-5} \def\xmax{5} % max min trục số
			\def\a{-2} % x<a
			\def\c{{)}} % Khoảng đoạn
			\draw[->] (\xmin,0)--(\a,0)node[scale=1.2]{\c}node[below=2.5pt]{$m$}--(\xmax,0);
			\begin{scope}[decoration={markings,mark=between positions 0 and 1 step 1mm
					with {\draw[blue!50] (70:5pt)--(250:5pt);}}]
				\draw[postaction={decorate}] (\a+0.1,0) -- (\xmax-0.1,0);
			\end{scope}
		\path (current bounding box.south) node[below=2mm] {$A=(-\infty ; m)$};
		\end{tikzpicture}
		\begin{tikzpicture}[line join = round, line cap = round,>=stealth,font=\footnotesize,scale=1]
		 \def\xmin{-5} \def\xmax{5} % max min trục số
		 \def\a{-3} \def\b{3} % a<x<b
		 \def\c{{[}} \def\d{{]}} % Khoảng đoạn
		 \draw[->] (\xmin,0)--(\a,0)node[scale=1.2]{\c}node[below=2.5pt]{$3m-1$}--(\b,0)node[scale=1.2]{\d}node[below=2.5pt]{$3m+3$}--(\xmax,0);
 		\begin{scope}[decoration={markings,mark=between positions 0 and 1 step 1mm
 		with {\draw[blue!50] (70:5pt)--(250:5pt);}}]
 		\draw[postaction={decorate}] (\xmin,0) -- (\a-0.1,0) (\b,0)--(\xmax-0.1,0);
 		\end{scope}
 		\path (current bounding box.south) node[below=2mm] {$B=[3m-1;3m+3]$};
		\end{tikzpicture}
		\end{center}
		\begin{itemchoice}
			\itemch \textbf{Sai}\\
			Khi $m=1$ thì tập $B=[2;6]$ chứa $5$ giá trị nguyên.
			\itemch \textbf{Đúng}\\
			Với $m=2$ thì $A=(-\infty;2)$ và $B=[5;9]$, suy ra $A \cap B=\varnothing$.
			\itemch \textbf{Sai}\\
			$B \subset A$ khi $3m+3<m \Leftrightarrow m<-\dfrac{3}{2}$.
			\itemch \textbf{Đúng}\\
			$C_{\mathbb{R}}A=[m;+\infty)$. Ta có $C_{\mathbb{R}}A \cap B \neq \varnothing$ khi $3m+3\ge m \Leftrightarrow m \geq-\dfrac{3}{2}$.
		\end{itemchoice}
	}
\end{ex}

\begin{ex}%[0-HK1-KN-4-ThanhBinh-DongNa-2324]%[VN-MT-25, Phạm Ngọc Trung]%Câu 4%[0D6H1-4]
	Một công ty sử dụng dây chuyền $A$ để đóng vào bao với khối lượng mong muốn là $5$ kg. Trên bao bì ghi thông tin khối lượng là $5 \pm 0{,}2$ kg. Gọi $\bar{a}$ là khối lượng thực của một bao gạo do dây chuyền $A$ đóng gói. 
	\choiceTF
	{Số đúng là $a=0{,}2$}
	{Số gần đúng là $\bar{a}=5{,}2$}
	{\True Độ chính xác là $d=0{,}2$}
	{\True Giá trị của $\bar{a}$ nằm trong đoạn $[4{,}8 ; 5{,}2]$}
	\loigiai{
		\begin{itemchoice}
			\itemch \textbf{Sai}\\
			Theo định nghĩa số đúng là $a=5$. Suy ra mệnh đề sai. 
			\itemch \textbf{Sai}\\
			Việc cho rằng $\bar{a} = 5{,}2$ là giá trị cao nhất có thể trong khoảng cho phép. Giá trị $\bar{a}$ hợp lý nhất trong trường hợp này phải là $5$ kg, vì nó chính là giá trị mà người sản xuất dự định làm giá trị trung tâm của khoảng dao động cho phép. 
			\itemch \textbf{Đúng}\\
			Độ chính xác là $d=0{,}2$. Suy ra mệnh đề đúng.
			\itemch \textbf{Đúng}\\
			Giá trị của $\bar{a}$ nằm trong đoạn $[4{,}8 ; 5{,}2]$.
		\end{itemchoice}
	}
\end{ex}

\Closesolutionfile{ans}
%================================PHẦN III
\caukq
\Opensolutionfile{ans}[ans/0-HK1-KN-4-ThanhBinh-DongNai-2324-Phan-III]

\begin{ex}%[0-HK1-KN-4-ThanhBinh-DongNa-2324]%[VN-MT-25, Phạm Ngọc Trung]%[0H5H4-1]
	Cho tam giác đều $ABC$ có cạnh bằng $4$. Giá trị của $\overrightarrow{BC} \cdot \overrightarrow{AB}$ bằng
	\shortans[]{$-4$}
	\loigiai{
	Ta có $\cos (\overrightarrow{BC}, \overrightarrow{AB})=- \cos (\overrightarrow{BC}, \overrightarrow{BA})=- \cos \widehat{ABC}=-\cos 60^{\circ}$.\\
	Do đó $\overrightarrow{BC} \cdot \overrightarrow{AB}=BC\cdot AB\cdot \cos (\overrightarrow{BC}, \overrightarrow{AB})=-4\cdot 4\cos 60^{\circ}=-4$.}
\end{ex}

\begin{ex}%[0-HK1-KN-4-ThanhBinh-DongNa-2324]%[VN-MT-25, Phạm Ngọc Trung]%[0D6H1-3]
	Biết điểm kiểm tra cuối học kỳ một của môn toán của một học sinh lớp $10$ là $\bar{a}=7{,}75$. Sai số tuyệt đối khi quy tròn điểm của học sinh này đến hàng phần chục là bao nhiêu?
	\shortans[]{$0{,}05$}
	\loigiai{Quy tròn điểm của học sinh này đến hàng phần chục ta được kết quả là $7{,}8$.\\
		Suy ra $\Delta_a=|7{,}75-7{,}8|=0{,}05$.}
\end{ex}

\begin{ex}%[0-HK1-KN-4-ThanhBinh-DongNa-2324]%[VN-MT-25, Phạm Ngọc Trung]%[1D5V1-4]
	Điểm kiểm tra thường xuyên môn toán của $40$ học sinh lớp $10$ được thống kê bằng bảng số liệu dưới đây (với $n$ là số nguyên dương và $n>3$). Nếu lớp có $40$ học sinh thì mốt của bảng số liệu thống kê đã cho là bao nhiêu?
	\begin{center}
		\begin{tabular}{|m{2cm}|m{1cm}|m{1cm}|m{1cm}|m{1cm}|m{1cm}|m{1cm}|m{1cm}|m{1cm}|}
			\hline
			Điểm & 3& 4& 5& 6& 7& 8& 9& 10\\
			\hline
			Số học sinh & 1& 2& $n+4$ & 4& 3& $3n-7$ & 4& 5\\
			\hline
		\end{tabular}
	\end{center}
\shortans[]{$8$}
	\loigiai{
	Theo giả thiết số học sinh bằng $40$ nên \[1+2+(n+4)+4+3+(3n-7)+4+5=40\Leftrightarrow 4n=24\Leftrightarrow n=6.\]
	Suy ra $n+4=10$, $3n-7=11$.\\
	Số học sinh đạt điểm $8$ nhiều nhất trong bảng số liệu trên nên $8$ là mốt của bảng số liệu.
	}
\end{ex}


\begin{ex}%[0-HK1-KN-4-ThanhBinh-DongNa-2324]%[VN-MT-25, Phạm Ngọc Trung]%[0D1V2-3]
Cho hai tập hợp $A=[-11; 15)$, $B=(2m-7;2m+1]$ khác rỗng. Tính tổng các giá trị nguyên của $m$ để $A\cup B=A$.
\shortans[]{$18$}
\loigiai{
\begin{center}
\begin{tikzpicture}[line join = round, line cap = round,>=stealth,font=\footnotesize,scale=1]
	\def\xmin{-5} \def\xmax{5} % max min trục số
	\def\a{-3} \def\b{3} % a<x<b
	\def\c{{[}} \def\d{{)}} % Khoảng đoạn
	\draw[->] (\xmin,0)--(\a,0)node[scale=1.2]{\c}node[below=2.5pt]{$-11$}--(\b,0)node[scale=1.2]{\d}node[below=2.5pt]{$15$}--(\xmax,0);
	\begin{scope}[decoration={markings,mark=between positions 0 and 1 step 1mm
			with {\draw[blue!50] (70:5pt)--(250:5pt);}}]
		\draw[postaction={decorate}] (\xmin,0) -- (\a-0.1,0) (\b,0)--(\xmax-0.1,0);
	\end{scope}
	\path (current bounding box.south) node[below=2mm] {$A=[-11; 15)$};
\end{tikzpicture}	
\begin{tikzpicture}[line join = round, line cap = round,>=stealth,font=\footnotesize,scale=1]
	\def\xmin{-5} \def\xmax{5} % max min trục số
	\def\a{-3} \def\b{3} % a<x<b
	\def\c{{(}} \def\d{{]}} % Khoảng đoạn
	\draw[->] (\xmin,0)--(\a,0)node[scale=1.2]{\c}node[below=2.5pt]{$2m-7$}--(\b,0)node[scale=1.2]{\d}node[below=2.5pt]{$2m+1$}--(\xmax,0);
	\begin{scope}[decoration={markings,mark=between positions 0 and 1 step 1mm
			with {\draw[blue!50] (70:5pt)--(250:5pt);}}]
		\draw[postaction={decorate}] (\xmin,0) -- (\a-0.1,0) (\b,0)--(\xmax-0.1,0);
	\end{scope}
	\path (current bounding box.south) node[below=2mm] {$B=(2m-7;2m+1]$};
\end{tikzpicture}	
\end{center}
Ta có $A\cup B=A\Leftrightarrow B\subset A\Leftrightarrow\heva{&2m-7\geq-11\\&2m+1< 15} \Leftrightarrow \heva{&m\geq-2\\&m < 7.}$\\
Khi đó tổng các giá trị nguyện $-2$, $-1$, $0$, $1$, $2$, $3$, $4$, $5$, $6$ bằng $18$.}
\end{ex}

\begin{ex}%[0-HK1-KN-4-ThanhBinh-DongNa-2324]%[VN-MT-25, Phạm Ngọc Trung]%[0H9V1-4]
Trong mặt phẳng tọa độ $Oxy$, cho hình bình hành $ABCD$ có $A(3;5)$, $B(7;2)$ và điểm $C$ thuộc trục hoành, điểm $D$ thuộc trục tung. Biết giao điểm $I$ của hai đường chéo của hình bình hành $ABCD$ có tọa độ là $(m; n)$. Tính giá trị của biểu thức $S=m+3n$.
\shortans[]{$11$}
\loigiai{
Mà điểm $C$ thuộc trục hoành nên $y_C=0$, điểm $D$ thuộc trục tung nên $x_D=0$.\\
Vì $I$ là giao điểm của hai đường chéo của hình bình hành $ABCD$ nên $I$ là trung điểm của $AC$ và $BD$. Nên ta có
\[ \heva{&2x_I=x_B+x_D\\&2y_I=y_A+y_C} \Rightarrow \heva{&2m=7\\&2n=5} \Leftrightarrow \heva{&m=\dfrac{7}{2}\\&n=\dfrac{5}{2}} \Rightarrow S=m+3n=\dfrac{7}{2}+\dfrac{15}{2}=11.\]}
\end{ex}

\begin{ex}%[0-HK1-KN-4-ThanhBinh-DongNa-2324]%[VN-MT-25, Phạm Ngọc Trung]%[0D2V2-3]
	Người ta dự định dùng hai loại nguyên liệu để chiết xuất ít nhất $140$ kg chất A và $9$ kg chất B. Từ mỗi tấn nguyên liệu loại I giá $4$ triệu đồng, có thể chiết xuất được $20$ kg chất A và $0{,}6$ kg chất B. Từ mỗi tấn nguyên liệu loại II giá $3{,}5$ triệu đồng, có thể chiết xuất được $10$ kg chất A và $1{,}5$ kg chất B. Biết rằng cơ sở cung cấp nguyên vật liệu chỉ có thể cung cấp không quá $10$ tấn nguyên liệu loại I và không quá $9$ tấn nguyên liệu loại II. Giả sử ta dùng $x$ (tấn) nguyên liệu loại I và $y$ (tấn) nguyên liệu loại II để chiết xuất các chất A và B. Chi phí mua nguyên vật liệu ít nhất là bao nhiêu? (\textit{đơn vị triệu đồng})
	\shortans[]{$34$}
	\loigiai{Ta dùng $x$ (tấn) nguyên liệu loại $I$ và $y$ (tấn) nguyên liệu loại I I để chiết xuất các chất A và B thỏa mãn yêu cầu đề bài.\\
	Ta có $\heva{& 0 \leq x \leq 10 \\ & 0 \leq y \leq 9}$.
	Tổng số tiền mua nguyên vật liệu là $T(x; y)=4 x+3{,}5 y$ triệu đồng.\\
	Số kg chất A chiết xuấ được là $20 x+10 y \Rightarrow 20 x+10 y \geq 140$.\\	
	Số kg chất B chiết xuất được là $0{,}6 x+1{,}5 y \Rightarrow 0{,}6 x+1{,}5 y \geq 9$.	
	Bài toán đã cho trở thành.\\
	Tìm các số $x$ và $y$ thỏa mãn hệ bất phương trình\allowdisplaybreaks
	\begin{eqnarray*}
		\heva{& 140 \leq 20 x+10 y \\ & 9 \leq 0{,}6 x+1{,}5 y\\ & 0 \leq x \leq 10 \\ & 0 \leq y \leq 9 } \Leftrightarrow\heva{& 14 \leq 2 x+y \\ & 30 \leq 2 x+5 y\\&0 \leq x \leq 10 \\ & 0 \leq y \leq 9 }
	\end{eqnarray*} sao cho $T(x; y)=4 x+3{,}5 y$ có giá trị nhỏ nhất.
	\begin{center}
		\begin{tikzpicture}[line join=round, line cap=round,>=stealth,thick,scale=.6]
		\tikzset{every node/.style={scale=0.9},
			declare function={x1=-1.5;x2=12.5;y1=-1.5;y2=11.5;}
		}
		\tikzset{declare function={
				a=-2;b=14;
				f(\x)=a*(\x)+b;
				a2=-2/5;b2=6;
				g(\x)=a2*(\x)+b2;	
				a3=-3/2;b3=27/2;
				h(\x)=a3*(\x)+b3;	
		}}
		\begin{scope}
		\clip (x1,y1)rectangle (x2-.3,y2-.3);
		\draw[smooth,samples=100,domain=x1:x2]plot(\x,{f(\x)});
		\fill[pattern=north west lines] plot[domain=x1:x2](\x,{f(\x)})--(x1,y1)--(x1,y2);
		
		\draw[smooth,samples=100,domain=x1:x2]plot(\x,{g(\x)});
		\fill[pattern=north west lines] plot[domain=x1:x2](\x,{g(\x)})--(x2,y1)--(x1,y1);
		
		%		\draw[smooth,samples=100,domain=x1:x2]plot(\x,{h(\x)});
		%		\fill[pattern=north west lines] plot[domain=x1:x2](\x,{h(\x)})--(x2,y2)--(x1,y2);
		%		
		\filldraw[pattern=north west lines] (x1-.1,9)rectangle(x2,y2);
		\filldraw[pattern=north west lines] (10,y1-.1)rectangle(x2-.1,y2+.1);
		\filldraw[pattern=north west lines] (x1-.1,y1-.1)rectangle(x2,0);
		\filldraw[pattern=north west lines] (x1-.1,y1-.1)rectangle(0,y2);
		\end{scope}
		\draw[->] (x1,0)--(x2,0)node[above]{$x$};
		\draw[->](0,y1)--(0,y2)node[left]{$y$};
		\filldraw[black] (0,0)circle(1.2pt)node[below left]{$O$};
		\foreach \x/\goc in {1,2,3,4,5,6,7,8,9,10}{
			\draw (\x,.15)--(\x,-.15)node[below]{$\x$};	
		}
		\draw 
		(5,4)circle(1pt)node[above right,scale=.8]{$A$} 
		(10,2)circle(1pt)node[above right,scale=.8]{$B$}
		(10,9)circle(1pt)node[above right,scale=.8]{$C$}
		(5/2,9)circle(1pt)node[above right,scale=.8]{$D$}		
		;
		\foreach \y/\goc in {1,2,3,4,5,6,7,8,9,10}{
			\draw (.15,\y)--(-.15,\y)node[left]{$\y$};	
		}
		\end{tikzpicture}
	\end{center}
	Ta xác định miền nghiệm $(S)$ của hệ $(*)$ là miền tứ giác $A B C D$ không được gạch chéo như hình vẽ. Trong đó $A(5; 4)$, $B(10; 2)$, $C(10; 9)$, $D\left(\dfrac{5}{2}; 9\right)$.\\
	Biểu thức $T(x; y)$ đạt giá trị nhỏ nhất tại một trong các đỉnh của tứ giác $ABCD$.\\
	Ta có $T(5{;}4)=34$, $T(10{;}2)=47$, $T(10{;}9)=71{,}5$, $T\left (\dfrac{5}{2}{;}9\right)=41{,}5$.\\
	Vậy chi phí mua nguyên vật liệu ít nhất bằng $T(5{,}4)=34$ triệu đồng.}
\end{ex}
\Closesolutionfile{ans}
% \begin{indapan}
% 	{ans/0-HK1-KN-4-ThanhBinh-DongNai-2324}
% \end{indapan}